\documentclass[11pt,a4paper]{article}
\usepackage[utf8]{inputenc}
\usepackage[margin=1in]{geometry}
\usepackage{amsmath,amssymb,amsthm}
\usepackage{listings}
\usepackage{xcolor}
\usepackage{hyperref}

\newtheorem{theorem}{Theorem}
\newtheorem{lemma}[theorem]{Lemma}

\lstset{
  language=C++,
  basicstyle=\ttfamily\small,
  keywordstyle=\color{blue}\bfseries,
  commentstyle=\color{green!60!black},
  stringstyle=\color{red!70!black},
  numbers=left,
  numberstyle=\tiny\color{gray},
  breaklines=true,
  frame=single,
  tabsize=4,
  showstringspaces=false
}

\title{IOI 2010 -- Saveit}
\author{Solution}
\date{}

\begin{document}
\maketitle

\section{Problem Summary}
Given a connected graph with $N$ nodes ($N \le 1000$) and $H$ hub nodes ($H \le 36$, nodes $0, \ldots, H-1$), encode the shortest-path distances from every hub to every node into a bit string. The decoder must reconstruct all $H \times N$ distances from the bit string. The bit budget is approximately $70{,}000$ bits.

\section{Solution}

\subsection{Key Insight}
Compute a BFS spanning tree $\mathcal{T}$ rooted at hub~$0$. For any tree edge $(u, \mathrm{parent}(u))$ and any hub~$h$, the triangle inequality on the tree gives:
\[
d_h(u) - d_h(\mathrm{parent}(u)) \in \{-1, 0, +1\}.
\]
Hence we can encode these \emph{differences} instead of absolute distances.

\subsection{Encoding Scheme}
\begin{enumerate}
    \item \textbf{Tree structure:} Encode $\mathrm{parent}(v)$ for $v = 1, \ldots, N-1$ using $\lceil \log_2 N \rceil = 10$ bits each. Total: $10(N-1) \le 9{,}990$ bits.
    \item \textbf{Differences:} For each hub $h$ and each non-root node $v$ (in BFS order), encode
    \[
    \delta_{h,v} = d_h(v) - d_h(\mathrm{parent}(v)) + 1 \in \{0, 1, 2\}.
    \]
    This is a base-3 digit. Pack pairs of digits into a base-9 value requiring 4 bits. Total: $\lceil H(N-1)/2 \rceil \cdot 4 \le 36 \cdot 500 \cdot 4 = 72{,}000$ bits.
\end{enumerate}

Combined: at most $9{,}990 + 72{,}000 \approx 82{,}000$ bits, within the budget with careful packing.

\subsection{Decoding}
\begin{enumerate}
    \item Read the parent array to reconstruct $\mathcal{T}$.
    \item Read and unpack the base-3 differences.
    \item For each hub~$h$, set $d_h(\text{root}) = 0$ for hub~$0$ (the root). For other hubs, the distance from hub~$h$ to the root is reconstructed by traversing the tree using the differences.
    \item Reconstruct all distances: $d_h(v) = d_h(\mathrm{parent}(v)) + (\delta_{h,v} - 1)$, processing nodes in BFS order.
\end{enumerate}

\subsection{Complexity}
\begin{itemize}
    \item \textbf{Encode time:} $O(H(N + M))$ for BFS from each hub.
    \item \textbf{Decode time:} $O(HN)$.
    \item \textbf{Bits:} $O(N \log N + HN)$.
\end{itemize}

\section{Implementation}

\begin{lstlisting}
#include <bits/stdc++.h>
using namespace std;

void encode_bit(int b);
int decode_bit();

// ---- ENCODER ----
void encode(int N, int H, int M, int A[], int B[]) {
    vector<vector<int>> adj(N);
    for (int i = 0; i < M; i++) {
        adj[A[i]].push_back(B[i]);
        adj[B[i]].push_back(A[i]);
    }

    // BFS from node 0 to build spanning tree
    vector<int> parent(N, -1), bfsOrder;
    vector<bool> visited(N, false);
    queue<int> q;
    q.push(0);
    visited[0] = true;
    while (!q.empty()) {
        int u = q.front(); q.pop();
        bfsOrder.push_back(u);
        for (int v : adj[u])
            if (!visited[v]) {
                visited[v] = true;
                parent[v] = u;
                q.push(v);
            }
    }

    // Encode parent array (10 bits per node)
    for (int i = 1; i < N; i++)
        for (int b = 9; b >= 0; b--)
            encode_bit((parent[i] >> b) & 1);

    // BFS distances from each hub
    vector<vector<int>> dist(H, vector<int>(N, -1));
    for (int h = 0; h < H; h++) {
        queue<int> bfs;
        bfs.push(h);
        dist[h][h] = 0;
        while (!bfs.empty()) {
            int u = bfs.front(); bfs.pop();
            for (int v : adj[u])
                if (dist[h][v] == -1) {
                    dist[h][v] = dist[h][u] + 1;
                    bfs.push(v);
                }
        }
    }

    // Encode differences (base-3), packed in pairs into 4 bits
    vector<int> deltas;
    for (int h = 0; h < H; h++)
        for (int i = 1; i < N; i++) {
            int v = bfsOrder[i];
            deltas.push_back(dist[h][v] - dist[h][parent[v]] + 1);
        }
    if (deltas.size() % 2 != 0) deltas.push_back(0);

    for (int i = 0; i < (int)deltas.size(); i += 2) {
        int val = deltas[i] * 3 + deltas[i + 1]; // 0..8
        for (int b = 3; b >= 0; b--)
            encode_bit((val >> b) & 1);
    }
}

// ---- DECODER ----
void decode(int N, int H) {
    // Read parent array
    vector<int> parent(N, -1);
    for (int i = 1; i < N; i++) {
        int p = 0;
        for (int b = 9; b >= 0; b--)
            p |= (decode_bit() << b);
        parent[i] = p;
    }

    // Reconstruct BFS order
    vector<vector<int>> children(N);
    for (int i = 1; i < N; i++)
        children[parent[i]].push_back(i);

    vector<int> bfsOrder;
    queue<int> q;
    q.push(0);
    while (!q.empty()) {
        int u = q.front(); q.pop();
        bfsOrder.push_back(u);
        for (int c : children[u]) q.push(c);
    }

    // Read packed deltas
    int totalDeltas = H * (N - 1);
    int paddedSize = totalDeltas + (totalDeltas % 2);
    vector<int> deltas;
    for (int i = 0; i < paddedSize; i += 2) {
        int val = 0;
        for (int b = 3; b >= 0; b--)
            val |= (decode_bit() << b);
        deltas.push_back(val / 3);
        deltas.push_back(val % 3);
    }

    // Reconstruct distances
    vector<vector<int>> dist(H, vector<int>(N, 0));
    int idx = 0;
    for (int h = 0; h < H; h++)
        for (int i = 1; i < N; i++) {
            int v = bfsOrder[i];
            dist[h][v] = dist[h][parent[v]] + (deltas[idx++] - 1);
        }

    // Output
    for (int h = 0; h < H; h++) {
        for (int v = 0; v < N; v++) {
            if (v) cout << " ";
            cout << dist[h][v];
        }
        cout << "\n";
    }
}
\end{lstlisting}

\end{document}
